\begin{itemize}
\item Účel: jasně definujte, zda asistent slouží jako tutor, generuje administrativní odpovědi, provádí předběžnou diagnostiku nebo facilitačně podporuje diskusi.
\item Kompetence: vysvětlovat koncepty v různé hloubce, navrhovat cvičení, kontrolovat formální náležitosti a indikovat možné chyby.
\item Hranice: nesmí poskytovat kompletní řešení hodnocených úloh, měnit hodnocení bez lidské kontroly ani zpracovávat citlivá data bez schváleného workflow.
\item Eskalace: plná řešení směrovat učiteli; při podezření na podvody oznamovat kontaktnímu pracovníkovi/katedře.
\item Eskalace (pokrač.): technické chyby či opakované nesprávnosti hlásit správci a rollbackovat konfiguraci.
\item Výstupy pro studenty: krátký hint (1–3 věty) a diagnostická poznámka „pravděpodobná oblast chyby“ (viz "Pokrok v matematice") \cite{kb_ai_101_tipu_pdf_04e4a115_page_8_1}.
\item Dále: výzva k nahrání mezikroků jako podmínka ověření postupu \cite{kb_ai_101_tipu_pdf_04e4a115_page_8_1}.
\item Výstupy pro učitele: agregovaný přehled problémových oblastí třídy, ukázky typických chyb a log konverzací/configurací (v souladu s GDPR).
\item Dokumentace a školení: doporučená nastavení a trénink pro pedagogy — např. Microsoft Learn \cite{kb_ai_101_tipu_pdf_04e4a115_page_54_1}.
\item Rychlý checklist před spuštěním: definujte role a hranice, zaveďte pravidla eskalace a kontakty, určete kdy požadovat nahrání postupu.
\item Připravte krátké školení pro pedagogy a komunikaci pro studenty o očekáváních a deklaraci využití AI.
\end{itemize}